% ===========================================================================
%  utente/05-projects/00-capitolo.tex
% ===========================================================================
\chapter{Projects}\label{cap:projects}

Un \voce{project} è una combinazione di più strategie, con il criterio che
decide quanto pesa ciascuna e gli eventuali freni. È ciò che si manda in
produzione: raramente conviene far girare una strategia sola.

\figuraSchermata{1.0}{figure/schermate/projects.png}
  {L'elenco dei project salvati.}
  {projects}

\section{Il costruttore, in quattro passi}\label{sez:costruttore}

Si apre con \voce{Ensemble builder}, oppure con \pulsantesec{Use in
Ensemble} dalla scheda di uno stream. I quattro passi in alto dicono sempre
a che punto si è.

\subsection{Passo 1 --- Strategies}

Si scelgono gli stream da combinare. Ne servono almeno due: con uno solo non
c'è niente da combinare, e il pulsante \pulsante{Next} resta spento finché
non se ne sono scelti due.

\figuraSchermata{1.0}{figure/schermate/composizione.png}
  {La scelta degli stream. Accanto a ciascuno ci sono Sharpe e rendimento,
   per non doverli cercare altrove.}
  {composizione}

\begin{perche}{Non scegliere i tre migliori}
La tentazione è prendere i tre Sharpe più alti. Quasi sempre è l'errore: se
le tre strategie fanno la stessa scommessa, metterle insieme non riduce
niente --- la si sta solo facendo in tre modi. La colonna \voce{Corr} nella
schermata iniziale serve a questo: due strategie con correlazione bassa, pur
con Sharpe modesti, insieme fanno meglio di tre strategie brave e uguali.
\end{perche}

\subsection{Passo 2 --- Weighting}

Come distribuire il capitale fra gli stream scelti. Quattordici metodi,
raggruppati per famiglia.

\figuraSchermata{1.0}{figure/schermate/composizione-pesi.png}
  {I metodi di pesatura, per famiglia. Ciascuno porta la sua riga di
   spiegazione.}
  {composizione-pesi}

\begin{center}
\begin{tabular}{@{}L{36mm}L{102mm}@{}}
\toprule
\textbf{Famiglia} & \textbf{L'idea} \\
\midrule
\rigaTenue
\voce{Naive} & Tutti uguali. \voce{Equal}, un mezzo a testa. Da provare
  sempre: è il riferimento contro cui misurare gli altri. \\
\voce{Risk / covariance} & Pesa in base a quanto ciascuno oscilla e a quanto
  si muovono insieme. Chi oscilla meno pesa di più. \\
\rigaTenue
\voce{Robust} & Come sopra, ma senza fidarsi troppo delle correlazioni
  stimate --- che su pochi dati sono instabili. \\
\voce{Tail-risk} & Guarda le perdite estreme invece dell'oscillazione
  media. \\
\rigaTenue
\voce{Return-aware} & Tiene conto anche di quanto ci si aspetta che ciascuno
  renda, non solo di quanto rischia. \\
\voce{Regime-adaptive} & Cambia i pesi secondo lo stato del mercato. \\
\bottomrule
\end{tabular}
\end{center}

\begin{pratica}{Da dove partire}
\voce{Equal}. Poi si prova \voce{Min-Variance LW} o \voce{ERC} e si guarda
se il drawdown scende davvero. Se un metodo complicato non batte l'equamente
pesato, l'equamente pesato è la risposta: ha un parametro in meno da
sbagliare.
\end{pratica}

\subsection{Passo 3 --- Overlay}

Gli overlay riducono l'esposizione quando succede qualcosa. Si possono
attivare più overlay insieme: i loro effetti si moltiplicano, e l'ordine in
cui li si sceglie non cambia il risultato.

\begin{center}
\begin{tabular}{@{}L{34mm}L{104mm}@{}}
\toprule
\textbf{Overlay} & \textbf{Riduce l'esposizione quando\ldots} \\
\midrule
\rigaTenue
\voce{Vol-Regime} & la volatilità realizzata sale. \\
\voce{Jump-Vol} & c'è un salto improvviso: reagisce in un giorno. \\
\rigaTenue
\voce{Corr-Gate} & le strategie cominciano a muoversi tutte insieme, cioè
  quando la diversificazione sparisce. \\
\voce{Regime-GMM}, \voce{Iso-Stress} & le condizioni di mercato somigliano a
  quelle dei periodi di stress passati. \\
\rigaTenue
\voce{Macro-Regime} & i segnali macroeconomici peggiorano. \\
\voce{News-Pressure} & il volume di notizie supera la norma. \\
\bottomrule
\end{tabular}
\end{center}

\begin{attenzione}{Un overlay toglie anche i guadagni}
Un freno che riduce l'esposizione riduce tutto: le perdite e i guadagni.
Vale la pena guardare i due numeri insieme --- di quanto è sceso il
drawdown, e di quanto è sceso il rendimento. Attivarne cinque perché «male
non fa» porta a un portafoglio quasi sempre fuori mercato.
\end{attenzione}

\subsection{Passo 4 --- Costs}

I costi: commissioni e slippage. Il backtest li sottrae.

\begin{nota}{Alzare i costi è il controllo più onesto che si possa fare}
Una strategia che resta buona raddoppiando i costi è robusta. Una che
smette di funzionare stava guadagnando dentro lo spread, e su un conto vero
non guadagnerà.
\end{nota}

\section{Salvare}\label{sez:projects-salvare}

Alla fine si dà un nome e si salva. Da quel momento il project compare
nell'elenco, il suo andamento è quello che la schermata iniziale disegna,
e --- cosa più importante --- si può scegliere nella pagina \voce{Live}
(\vediCap{live}).
